\IfFontExistsTF{Noto Serif CJK SC}{
  \setCJKmainfont{Noto Serif CJK SC}[BoldFont={Noto Serif CJK SC Bold},AutoFakeSlant=0.15]
  \setCJKsansfont{Noto Sans CJK SC}
  \setCJKmonofont{Noto Sans Mono CJK SC}
}{
  \setCJKmainfont{FandolSong-Regular.otf}[BoldFont=FandolSong-Bold.otf,ItalicFont=FandolKai-Regular.otf]
  \setCJKsansfont{FandolHei-Regular.otf}[BoldFont=FandolHei-Bold.otf]
  \setCJKmonofont{FandolFang-Regular.otf}
}
\usepackage[a4paper,margin=25mm,headheight=15pt]{geometry}
\usepackage{amsmath,amssymb,amsthm,bm,mathtools}
\usepackage{graphicx,booktabs,longtable,array,tabularx,multirow}
\usepackage{enumitem,xcolor,caption,float}
\usepackage[unicode,hypertexnames=false,colorlinks=true,linkcolor=blue!45!black,citecolor=green!35!black,urlcolor=blue!55!black]{hyperref}
\usepackage{bookmark,fancyhdr}
\setlist{topsep=5pt,leftmargin=2em}
\setcounter{tocdepth}{1}
\setlength{\parskip}{3pt}
\setlength{\emergencystretch}{2em}
\newcommand{\R}{\mathbb R}
\newcommand{\C}{\mathbb C}
\newcommand{\N}{\mathbb N}
\newcommand{\Z}{\mathbb Z}
\newcommand{\E}{\mathbb E}
\newcommand{\Prob}{\mathbb P}
\newcommand{\T}{\mathsf T}
\newcommand{\HH}{\mathsf H}
\newcommand{\I}{\mathbf I}
\newcommand{\one}{\mathbf 1}
\newcommand{\norm}[1]{\left\lVert #1\right\rVert}
\newcommand{\abs}[1]{\left\lvert #1\right\rvert}
\newcommand{\inner}[2]{\left\langle #1,#2\right\rangle}
\DeclareMathOperator{\dom}{dom}
\DeclareMathOperator{\tr}{tr}
\DeclareMathOperator{\diag}{diag}
\DeclareMathOperator{\rank}{rank}
\DeclareMathOperator{\cl}{cl}
\DeclareMathOperator{\relint}{relint}
\DeclareMathOperator{\conv}{conv}
\DeclareMathOperator{\aff}{aff}
\DeclareMathOperator{\dist}{dist}
\DeclareMathOperator{\vol}{vol}
\DeclareMathOperator{\epi}{epi}
\DeclareMathOperator*{\argmin}{arg\,min}
\DeclareMathOperator*{\argmax}{arg\,max}
\newenvironment{origexample}[1]{\par\medskip\noindent\textbf{例 #1}\quad\ignorespaces}{\par\medskip}
\newenvironment{origalgorithm}[1]{\par\medskip\noindent\textbf{算法 #1}\quad\ignorespaces}{\par\medskip}
\pagestyle{fancy}\fancyhf{}
\fancyhead[L]{\small 凸优化}
\fancyhead[R]{\small\nouppercase{\leftmark}}
\fancyfoot[C]{\thepage}
\ctexset{chapter={format=\huge\bfseries,beforeskip=8pt,afterskip=22pt},section={format=\Large\bfseries}}
\numberwithin{equation}{chapter}
\graphicspath{{figures/}}
\hypersetup{pdftitle={凸优化：中文译稿},pdfauthor={Stephen Boyd; Lieven Vandenberghe}}
